\chapter{Solution Implementation}

\section{Project Structure}
This chapter details how the architecture from the previous chapter is realized in
the running \textit{Dataspecer Adjuster} implementation. Where relevant, it
cross-references decisions introduced in the Solution Design and explains concrete
module boundaries, runtime wiring, and operational aspects.

\subsection*{Repository layout}
The Adjuster is implemented as a multi-container, monorepo-style workspace with a
NestJS monorepo for the backend and a Next.js application for the review UI:
\begin{itemize}
  \item \texttt{services/adjuster/backend}: NestJS monorepo with apps
  \texttt{dialogs-handler}, \texttt{changes-detector}, \texttt{changes-suggester},
  \texttt{dataspecer-adapter}, and \texttt{artifacts-embedder}, plus a shared library
  \texttt{libs/common} for DTOs, prompts, utilities and a minimal MCP shim.
  \item \texttt{services/adjuster/frontend/dataspecer-validator-ui}: Next.js (React) UI for
  the two-panel diff review, powered by Tailwind and \texttt{jsondiffpatch}.
  \item \texttt{services/adjuster/docker-compose.yml}: development orchestration of backend
  apps, the UI, PostgreSQL, RabbitMQ, Qdrant, and a local Dataspecer instance.
  \item \texttt{services/adjuster/docker-compose.prod.yml}: production-oriented compose file
  with Nginx reverse proxy and hardened env wiring.
  \item \texttt{services/adjuster/reverseproxy}: Nginx configuration for TLS termination and routing.
\end{itemize}

The backend is built with NestJS~11 using the monorepo setup (\texttt{nest-cli.json})
that hosts independent applications (``containers'' in the C4 sense) while enabling
shared types and utilities via the common library. The UI is a separate Next.js~15 app.

\subsection*{Backend applications}
Each backend app encapsulates a responsibility from the design:
\begin{itemize}
  \item \textbf{Dialogs Handler} (\texttt{apps/dialogs-handler}): Orchestrates sessions and exposes the HTTP API for developer and maintainer flows (listing specifications, uploading schemas, running analyses, applying accepted changes, exporting reports, generating validation links, and LLM chat endpoints). It persists conversational and evaluation artifacts to PostgreSQL via Prisma.
  \item \textbf{Changes Detector} (\texttt{apps/changes-detector}): Performs deterministic ingestion and differencing of JSON Schemas with a structured LLM-backed classifier. It can fetch the baseline JSON Schema or PSM via the \texttt{dataspecer-adapter}.
  \item \textbf{Changes Suggester} (\texttt{apps/changes-suggester}): Encapsulates the LLM prompting and structured output normalization for change rationales and suggestion phrasing, aligned with the output schemas in \texttt{libs/common/dto}.
  \item \textbf{Dataspecer Adapter} (\texttt{apps/dataspecer-adapter}): Provides a thin boundary to Dataspecer’s HTTP/MCP endpoints. It centralizes export/fetch operations and translates internal IRIs and payloads into Dataspecer API calls.
  \item \textbf{Artifacts Embedder} (\texttt{apps/artifacts-embedder}): Periodically extracts artifacts (PSM slices, JSON Schemas) and embeds them into a Qdrant vector store to support RAG in the detector/suggester flows.
\end{itemize}
The shared library \texttt{libs/common} provides:
\begin{itemize}
  \item DTOs for requests/responses across applications (diff requests, detected changes, evaluation metrics, chat payloads).
  \item Prompt templates for the LLM interaction and helper utilities for PSM context slicing.
  \item A minimal MCP server shim for typed tool exposure with authentication.
\end{itemize}

\subsection*{Frontend application}
The UI (\texttt{dataspecer-validator-ui}) is a Next.js~15 app. It implements a compact
review interface with a diff tree and a details/decision panel. Environment variables
bind it to backend HTTP endpoints (Dialogs Handler, Changes Detector/Suggester) and the
Dataspecer backend. The UI favors clarity and responsiveness over heavy visualization,
as motivated in the design chapter.

\section{Used Approaches}
This section explains how the core approaches outlined in the Solution Design were implemented:
structured differencing, LLM guardrails, adapter boundary, MCP tooling, and operations.

\subsection{NestJS modularity and DI}
The monorepo is configured via \texttt{nest-cli.json} with one project per application
and a shared \texttt{common} library. Each app has a root module and composes providers
through Nest’s dependency injection. Inter-app coordination happens via:
\begin{itemize}
  \item \textbf{HTTP} (Dialogs Handler public API for the UI),
  \item \textbf{AMQP (RabbitMQ)} for asynchronous flows and decoupled calls (e.g., retrieving JSON Schema/PSM via the adapter),
  \item \textbf{Shared DTOs} and contracts from \texttt{libs/common}.
\end{itemize}
This matches the container split from the Solution Design while keeping the codebase DRY and testable.

\subsection{Dependency injection and decorators (NestJS)}
The implementation relies on NestJS’s DI container and declarative decorators:
\begin{itemize}
  \item \textbf{Providers and constructor injection}: services are marked with \texttt{@Injectable()} and injected via constructors. Where an explicit token is needed (e.g., messaging clients), \texttt{@Inject('DATASPECER\_ADAPTER')} is used to resolve a \texttt{ClientProxy} for AMQP calls.
  \item \textbf{HTTP controllers}: request handling is declared with \texttt{@Controller()}, method-level \texttt{@Get}, \texttt{@Post}, and parameter decorators \texttt{@Body}, \texttt{@Query}, \texttt{@Param}. Public API documentation is generated through Swagger decorators such as \texttt{@ApiTags}, \texttt{@ApiOperation}, \texttt{@ApiBody}, \texttt{@ApiOkResponse}.
  \item \textbf{Scheduling}: background refresh is implemented with \texttt{@Cron(...)} in the Artifacts Embedder, enabling periodic embedding of Dataspecer artifacts.
  \item \textbf{DTOs and validation}: request/response contracts live in DTO classes and are validated using \texttt{class-validator}. This keeps HTTP and AMQP boundaries type-safe and reject invalid payloads early.
  \item \textbf{Microservices}: AMQP-facing controllers and services communicate through Nest microservices APIs; DI supplies the messaging client and adapter boundary, decoupling detector/suggester from Dataspecer specifics.
\end{itemize}
This decorator- and DI-centric style concentrates routing, typing, and documentation in metadata while keeping business logic in injectable services, consistent with the modular split described in the Solution Design.

\subsection{Deterministic diff + LLM classification}
The \textbf{Changes Detector} accepts either a pair of JSON Schemas or a pair consisting of a new JSON Schema and a PSM/JSON fetched via \texttt{dataspecer-adapter}. In implementation:
\begin{itemize}
  \item Structural differences are computed using \texttt{json-diff}. For large inputs, the diff is truncated to bounded sizes to keep latency and cost predictable.
  \item Diff keywords are extracted and used to slice a relevant PSM context window (``local'' RAG with no external store dependency).
  \item A task-specific prompt from \texttt{libs/common/config/prompts.ts} is filled with the truncated diff and the PSM slice.
  \item Outputs are enforced via \textbf{Zod} schemas using LangChain’s structured output mode; only conformant outputs are accepted.
\end{itemize}
This realizes the \emph{detect} $\rightarrow$ \emph{classify} $\rightarrow$ \emph{explain} path from the design, preserving determinism and auditability.

\subsection{LLM scaffolding and safety}
LLM calls are made through \texttt{@langchain/openai} with conservative temperature and token caps. Guardrails:
\begin{itemize}
  \item \textbf{Strict output schemas} (Zod) for changes and rationales. Non-conforming generations are rejected.
  \item \textbf{Prompt discipline}: no raw specification content is echoed back to the user; only concise rationales and structured summaries are produced for supported change types.
  \item \textbf{Time/size bounds}: truncation of diffs and contexts, with sensible limits configured in the service.
\end{itemize}
Where the LLM path fails, the detector returns a transparent fallback item instructing manual review.

\subsection{Dataspecer adapter boundary}
The \textbf{Dataspecer Adapter} encapsulates all communication with Dataspecer services. The implementation provides:
\begin{itemize}
  \item Fetching JSON Schema and PSM by IRI for detector/suggester flows.
  \item A single place to evolve protocol details or authentication (e.g., switching between direct REST and MCP-backed endpoints).
  \item Idempotent patterns for apply operations (create/update properties, toggle \texttt{required}, update enums) with a dry-run-first discipline enforced by the Dialogs Handler before committing mutations.
\end{itemize}
This adapter boundary keeps Dataspecer evolution localized and simplifies change impact analysis.

\subsection{Minimal MCP server and tools}
To support tool-based orchestration (and enable external MCP-capable clients), the shared library exposes a minimal \textbf{MCP shim}:
\begin{itemize}
  \item An SSE endpoint emits \texttt{mcp.hello} with the catalog of available tools.
  \item A JSON-RPC-style message endpoint handles tool invocations; an authentication guard accepts either a bearer token or an HMAC signature over a date header.
  \item A \texttt{health} tool is implemented end-to-end; Dataspecer-oriented tools (\texttt{list\_specs}, \texttt{get\_psm}, \texttt{preview\_apply}, \texttt{apply\_changes}, \texttt{validate}) can register to this surface and reuse the same auth and audit hooks.
\end{itemize}
This realizes the design goal of a typed, preview-first tool layer and keeps the safety properties observable and auditable.

\subsection{Persistence and audit}
Dialog and evaluation artifacts are stored in PostgreSQL through \textbf{Prisma}. The current schema includes chat messages and specification maintainer sessions with analysis results, decisions, and status tracking. This enables:
\begin{itemize}
  \item Reproducible evaluation runs (diff quality, apply success).
  \item End-to-end auditability of decisions and MCP actions.
  \item Reopening and continuing sessions by both roles.
\end{itemize}
The presence of TypeORM in dependencies reflects initial exploration; the implementation standardizes on Prisma for dialog-centric workflows to align with typed client ergonomics and easier migrations.

\subsection{Dialogs Handler HTTP API}
The Dialogs Handler exposes REST endpoints covering the user stories from the Solution Design:
\begin{itemize}
  \item \textbf{Specifications}: list and fetch by IRI (mocked in development, adapter-backed in integration).
  \item \textbf{Analyses}: run checks against uploaded schemas; store and retrieve shared analyses.
  \item \textbf{Decisions}: accept/reject/flag with optional notes; generate developer reports.
  \item \textbf{Apply}: preview and apply changes (adapter-backed), enforcing a preview-before-apply discipline.
  \item \textbf{Validation links}: create tokenized links for developers to self-check re-uploads.
  \item \textbf{LLM chat}: start/send/get/delete conversations, regenerate change descriptions from maintainer feedback, and validate developer re-uploads against stored analyses.
\end{itemize}
All inputs/outputs are typed with DTOs and validated. CORS origins are controlled via environment variables.

\subsection{Artifacts Embedder and Qdrant}
The \textbf{Artifacts Embedder} extracts JSON Schemas and PSM fragments from Dataspecer, computes embeddings, and writes chunked payloads to Qdrant. It offers:
\begin{itemize}
  \item \textbf{Chunking}: configurable chunk size with stable chunk IDs.
  \item \textbf{Metadata}: payloads carry IRIs and types (psm/json-schema) to filter later retrievals.
  \item \textbf{Search}: cosine-similarity \texttt{search} with optional filter constraints (\texttt{psmIri}, \texttt{dataSpecificationIri}).
  \item \textbf{Cron refresh}: an optional scheduled refresh for a configured specification.
\end{itemize}
This component enables scalable RAG without coupling detector/suggester runtime to Dataspecer availability.

\subsection{Frontend implementation}
The Next.js UI implements:
\begin{itemize}
  \item A diff tree (left pane) and a detail/decision panel (right pane), with color-coded categories (adds/removes/renames/type changes).
  \item Links from changes to local diff snippets; concise rationales generated by the LLM when available.
  \item Iterative re-upload flow for developers using validation links.
\end{itemize}
Endpoints are configured via \texttt{NEXT\_PUBLIC\_*} environment variables to avoid hard-coded addresses and facilitate reverse proxy routing in production.

\subsection{DevOps and runtime topology}
The Compose files realize the container view from the Solution Design:
\begin{itemize}
  \item \textbf{Development} (\texttt{docker-compose.yml}): binds the UI to local backend ports, boots PostgreSQL, RabbitMQ, Qdrant, and a local Dataspecer instance; enables hot-reload for backend apps; exposes MCP endpoints behind the Dialogs Handler when enabled.
  \item \textbf{Production} (\texttt{docker-compose.prod.yml}): builds optimized images for each backend app, runs the Next.js UI separately, and fronts the stack with Nginx (\texttt{reverseproxy/nginx.conf}). Health checks and strict environment variables are used for readiness.
  \item \textbf{Security}: MCP requires a shared secret (bearer or HMAC), CORS is constrained, and only the Nginx port is public in the production topology.
\end{itemize}
Environment variables configure origins, Dataspecer endpoints, OpenAI credentials, MCP base paths and tokens, and DB/AMQP/Qdrant connections. Volumes persist Postgres and Qdrant state.

\subsection{Testing and observability}
Each backend app provides a Jest e2e configuration; shared DTOs stabilize contracts across tests. Basic structured logging is present in LLM-bound paths (input sizes, truncation, and result counts) to monitor cost and correctness. These logs, together with Prisma records, support the evaluation metrics defined in the Solution Design (precision/recall per change type, apply success rate, MCP safety indicators).

\subsection{Security and guardrails}
In line with the rationale from the design chapter:
\begin{itemize}
  \item \textbf{Guarded mutations}: apply operations only after explicit preview; MCP tools mirror this via a preview-first discipline.
  \item \textbf{LLM boundaries}: Zod-enforced outputs; size and time caps; conservative temperature.
  \item \textbf{Auth}: MCP requires bearer/HMAC; UI endpoints require CORS-allowed origins; tokens scope validation links.
\end{itemize}
This deliberately constrains the LLM to a recommender role and keeps mutations typed, previewed, and auditable.

\subsection{Limitations and future work}
Current limitations match the incremental scope:
\begin{itemize}
  \item The diff/classifier focuses on JSON Schema-level changes most common in practice; deeper OWL-level reasoning is out of scope but can be added behind the same adapter and MCP signatures.
  \item The minimal MCP shim currently exposes a health tool and a typed surface; a full Dataspecer toolset is straightforward to register and will enhance external orchestration.
  \item Retrieval in the detector is local-context-first; the Qdrant-backed embedder enables richer cross-artifact retrieval but is not yet integrated into all flows by default.
\end{itemize}
These choices align with the KISS/DRY principles and the evaluation goals, while keeping a clear path to richer semantics and expanded tool exposure.
